
\subsection{Sensitivity to the zone aggregation}
\label{subsec:clustering}

The zone-aggregated level rests on one hand-made seven-zone partition of the fifteen-node
network. A partition is a modelling choice rather than a physical fact, and a different one
could in principle redistribute losses and routing differently. The aggregated model was
therefore re-solved across a family of alternative partitions, holding everything else fixed.

The family comprises the original partition re-run through the same aggregator, three
deliberate alternatives spanning granularity (four, seven and ten zones) and a null
distribution of twenty random tree-contiguous partitions, all solved to a $0.01$\,\%
optimality gap. Every member conserves the total heat-loss capacity
$\sum_p U_p L_p = 1140.2$\,W/K of the node-resolved network, which is verified per
configuration and confirmed by an identical realised annual loss across the family
($1257.8$--$1257.9$\,MWh). Conservation requires resolving the producer as its own zone: a
partition that merges the plant with downstream nodes has no incoming pipe to carry the
producer zone's internal trunk losses and silently discards them. The producer is therefore held
isolated in every member, as in the original partition, and the consequence for the null is
stated explicitly: each draw cuts the producer edge and then
five further random edges, partitioning the remaining fourteen consumer nodes into six
zones. The null therefore tests consumer-zone boundary choice at fixed producer resolution
and fixed total loss, which is the comparison the question requires.

Because the delivered loss is equal across the family, the only quantity free to move is
where the boundaries fall, a pure routing effect. It barely moves. The three deliberate
alternatives differ from the original by at most $0.01$\,\% of annual cost, the twenty null
draws have a standard deviation of $4$\,EUR, and the full spread across all twenty-four
partitions is \textbf{11}\,EUR, or \textbf{0.008}\,\% of operating cost. Granularity between
four and ten zones changes essentially nothing.

Set against the decomposition, this is more than a robustness check. With losses held
equal, the residual is directly comparable to the spatial-topology main effect measured on
the same network, and the three quantities span three orders of magnitude
(Table~\ref{tab:r16}).

\begin{table}[htbp]
\centering
\footnotesize
\caption{Zone-boundary choice against the effects it must be compared with, on the
copperplate-to-baseline ($\CP\!\to\!\Lone$) cost gap. With the total heat-loss capacity held
equal across all twenty-four partitions, the residual routing effect is negligible.}
\label{tab:r16}
\begin{tabular}{lrr}
\hline
Effect & Cost & \% of the $\CP\!\to\!\Lone$ gap \\
\hline
Loss visibility & 19\,737\,EUR & 95.8 \\
Spatial topology (resolving routing, physics-matched) & 52\,EUR & 0.25 \\
Zone-boundary choice (this subsection) & $\lesssim$11\,EUR & at solver tolerance \\
\hline
\end{tabular}
\end{table}


Choosing where to draw a zone boundary is negligible: at $\lesssim$11\,EUR the residual is
three orders of magnitude below representing the loss and, like the physics-matched topology
effect itself, sits at the $0.01$\,\% solver tolerance
($\approx$13\,EUR on the Memmingen cost), so no zone-partition effect is resolved
above it. The aggregation is therefore not a hidden degree
of freedom in the results, and the finding corroborates the decomposition from an
independent direction: loss, not geometry, carries the cost.

The corollary is a practical warning worth stating. An aggregation that does \emph{not}
conserve the total heat-loss capacity (for instance one that merges the producer with
downstream nodes and so discards the trunk pipes internal to that zone, here $10$--$18$\,\%
of $\sum_p U_p L_p$) shifts reported cost by up to about $6$\,\%, against the $0.01$\,\%
agreement among conserving partitions. The binding requirement for a valid reduced network
model is preservation of the delivered loss, not fidelity of the zone geometry. That is the
same conclusion the decomposition reaches by a different route, and it is a straightforward
error to make.

One parameterisation note, since two values for the same pipe appear in the released data.
The aggregator uses the network's per-pipe design heat-transfer coefficients, under which
the far-east lateral carries a markedly higher value than the trunk, reflecting its older
uninsulated construction. The Stage-1 validation instead calibrates a single uniform
coefficient across all pipes, one free parameter, to reproduce the annual loss implied by the delivered-energy measurement.
The two serve different purposes and are not intended to agree pipe by pipe: the first is
the design network, the second a minimal calibration of its aggregate behaviour. All
partitions compared here use the same per-pipe basis, so the comparison is unaffected.
